\documentclass[aspectratio=169]{beamer}
\usetheme{Madrid}
\usecolortheme{whale}

\usepackage{graphicx}
\usepackage{booktabs}
\usepackage{hyperref}
\usepackage{multicol}
\usepackage{array}
\usepackage{xcolor}
\usepackage{colortbl}
\usepackage{tikz}

\definecolor{iitbblue}{RGB}{0,56,147}
\definecolor{lightgray}{RGB}{240,240,240}
\definecolor{darkgreen}{RGB}{0,128,0}
\definecolor{darkred}{RGB}{180,0,0}

\setbeamertemplate{navigation symbols}{}
\setbeamertemplate{footline}[frame number]

\title[EE 679: Movie Auto-Subtitling]{EE 679: Noise-Aware Adaptive Auto-Subtitling for Movies}
\subtitle{Classical Speech Processing + Whisper ASR + Scene Understanding}
\author{Aman Verma}
\institute{Department of Electrical Engineering\\Indian Institute of Technology Bombay}
\date{April 2026}

\begin{document}

% ─────────────────────────────────────────────────────────────────────────────
\begin{frame}
  \titlepage
\end{frame}

% ─────────────────────────────────────────────────────────────────────────────
\begin{frame}{Outline}
  \begin{columns}[T]
    \begin{column}{0.5\textwidth}
      \textbf{System Design}
      \begin{enumerate}
        \item Motivation \& Problem Statement
        \item Pipeline Architecture (6 Stages)
        \item Datasets \& Ground Truth
      \end{enumerate}
      \vspace{0.3cm}
      \textbf{Core Experiments (LibriSpeech)}
      \begin{enumerate}
        \setcounter{enumi}{3}
        \item VAD Ablation Study
        \item Enhancement Ablation
        \item Noise Robustness Benchmark
        \item Model Size: tiny.en vs medium.en
        \item Subtitle Quality \& Readability
      \end{enumerate}
    \end{column}
    \begin{column}{0.5\textwidth}
      \textbf{Trailer Evaluation}
      \begin{enumerate}
        \setcounter{enumi}{8}
        \item Mass Trailer Eval (14 trailers, 8 genres)
        \item The YouTube Ground Truth Fallacy
      \end{enumerate}
      \vspace{0.3cm}
      \textbf{Novel Experiments}
      \begin{enumerate}
        \setcounter{enumi}{10}
        \item PESQ \& STOI Enhancement Quality
        \item Noise-Type Classifier
        \item Whisper Confidence Calibration
        \item Streaming Pipeline Latency
      \end{enumerate}
      \vspace{0.3cm}
      \textbf{Synthesis}
      \begin{enumerate}
        \setcounter{enumi}{14}
        \item Scene Understanding \& Mood
        \item Key Findings \& Future Work
      \end{enumerate}
    \end{column}
  \end{columns}
\end{frame}

% ─────────────────────────────────────────────────────────────────────────────
\section{Motivation}
\begin{frame}{Motivation \& Problem Statement}
  \begin{columns}[T]
    \begin{column}{0.55\textwidth}
      \textbf{Why is Movie Subtitling Hard?}
      \begin{itemize}
        \item \textbf{Variable SNR}: dialogue at $+$20 dB, action sequences at $-$5 dB
        \item Dialogue overlaps with \textbf{orchestral score}, sound effects, crowd noise
        \item Standard Whisper \textbf{hallucinates} during non-speech segments (loops lyrics, repeats)
        \item \textbf{YouTube auto-captions} fail on cinematic content — they are flawed reference transcripts
        \item Accessibility requirement: subtitles must satisfy \textbf{$\leq$20 chars/sec} (CPS) for hearing-impaired viewers
      \end{itemize}
    \end{column}
    \begin{column}{0.42\textwidth}
      \begin{block}{Project Goal}
        Build a \textbf{noise-aware, end-to-end} auto-subtitling pipeline that:
        \begin{itemize}
          \item Transcribes accurately under movie-level noise
          \item Adapts enhancement to SNR conditions
          \item Annotates output with \textbf{scene boundaries, mood, and emotion}
          \item Operates \textbf{faster than real-time} on CPU
        \end{itemize}
      \end{block}
      \begin{alertblock}{Novelty}
        Adaptive routing that \emph{withholds} enhancement for clean segments — preventing over-processing that degrades Whisper.
      \end{alertblock}
    \end{column}
  \end{columns}
\end{frame}

% ─────────────────────────────────────────────────────────────────────────────
\section{Pipeline}
\begin{frame}{6-Stage Pipeline Architecture}
  \begin{columns}[T]
    \begin{column}{0.48\textwidth}
      \small
      \begin{block}{Stage 1: Audio Extraction}
        \texttt{ffmpeg} (bundled via imageio-ffmpeg)\\
        $\to$ mono 16 kHz PCM WAV
      \end{block}
      \begin{block}{Stage 2: Voice Activity Detection}
        Three methods compared:\\
        \textbf{(A)} Energy threshold (baseline)\\
        \textbf{(B)} MFCC C0 $\times$ ZCR\\
        \textbf{(C)} Full spectral (proposed):\\
        \quad $0.45\cdot C_0 + 0.25\cdot\text{centroid}$\\
        \quad $-0.20\cdot\text{flatness} + 0.10\cdot(1-\text{ZCR})$
      \end{block}
      \begin{block}{Stage 3: Adaptive Enhancement}
        Routing decision per segment:\\
        \textbf{IF} SNR $< 18$ dB \textbf{OR} flatness $> 0.18$\\
        \textbf{OR} ZCR $> 0.12$ $\to$ \textbf{enhance}\\
        \textbf{ELSE} $\to$ pass raw to ASR\\
        Options: Wiener filter / Spectral Subtraction
      \end{block}
    \end{column}
    \begin{column}{0.48\textwidth}
      \small
      \begin{block}{Stage 4: ASR — faster-whisper}
        \texttt{tiny.en} (39M) / \texttt{base.en} (74M) / \texttt{medium.en} (307M), int8\\
        Silero VAD: thresh=0.35, min\_silence=400 ms\\
        Hallucination suppression: log\_prob $< -1.2$\\
        Word-level timestamps for SRT alignment\\
        RTF: 0.035–0.23 (all $< 1$ = real-time)
      \end{block}
      \begin{block}{Stage 5: Scene Understanding}
        Foote novelty score (MFCC+chroma SSM)\\
        $\to$ scene boundaries\\
        Per-segment mood classification:\\
        \textit{tense / epic / sad / calm / dialogue / silence}
      \end{block}
      \begin{block}{Stage 6: Subtitle Post-Processing}
        $\leq$12 words/cue, max 84 chars/cue (2$\times$42)\\
        CPS enforcement ($>$20 = violation)\\
        Outputs: \texttt{.srt}, annotated SRT, JSON, PNG
      \end{block}
    \end{column}
  \end{columns}
\end{frame}

% ─────────────────────────────────────────────────────────────────────────────
\begin{frame}{Datasets \& Ground Truth}
  \begin{columns}[T]
    \begin{column}{0.48\textwidth}
      \textbf{LibriSpeech test-clean} (controlled benchmark)
      \begin{itemize}
        \item 2,620 utterances, 5.4 hours, 40 speakers
        \item 690 MB, auto-downloaded from openslr.org
        \item Used for: noise robustness, VAD, enhancement, model comparison
        \item \textbf{Synthetic noise added}: white Gaussian, pink (1/f), babble, movie soundtrack
        \item SNR range: $-5$ to $+20$ dB (5 dB steps)
      \end{itemize}
      \vspace{0.2cm}
      \textbf{Movie Trailers (YouTube)}
      \begin{itemize}
        \item \textbf{14 trailers} with reference SRTs, 8 genres
        \item Action, sci-fi, drama, comedy, horror,\\animation, thriller, biopic
        \item 246 MB (trailers) + 990 MB (youtube\_trailers)
        \item Reference: YouTube auto-captions (VTT$\to$SRT)
        \item \textbf{Official SRTs}: Avengers \& Spider-Man only
      \end{itemize}
    \end{column}
    \begin{column}{0.48\textwidth}
      \textbf{ASR Models Evaluated}
      \begin{center}
        \small
        \begin{tabular}{@{}lccc@{}}
          \toprule
          Model & Params & WER & RTF \\
          \midrule
          \texttt{tiny.en} & 39M & 0.875 & 0.052 \\
          \texttt{base.en} & 74M & 0.172 & 0.046 \\
          \texttt{medium.en} & 307M & \textbf{0.522} & 0.175 \\
          \bottomrule
        \end{tabular}
      \end{center}
      \vspace{0.3cm}
      \begin{alertblock}{The Ground Truth Problem}
        YouTube auto-captions hallucinate during music and non-speech segments. Our pipeline correctly outputs \textbf{silence}; the reference outputs lyrics. \\
        $\Rightarrow$ WER against official SRTs: \textbf{0.144} (Avengers)\\
        $\Rightarrow$ WER against YouTube SRTs: $\sim$0.52 (inflated)
      \end{alertblock}
    \end{column}
  \end{columns}
\end{frame}

% ─────────────────────────────────────────────────────────────────────────────
\section{Core Experiments}
\begin{frame}{Experiment 3: VAD Ablation Study}
  \begin{columns}[T]
    \begin{column}{0.55\textwidth}
      \textbf{Setup}: Compare 3 VAD methods on speech coverage.\\
      \textbf{Dataset}: LibriSpeech + Avengers trailer (music-heavy).

      \vspace{0.3cm}
      \begin{center}
        \begin{tabular}{@{}lccc@{}}
          \toprule
          \textbf{VAD Method} & \textbf{Coverage} & \textbf{Compute} & \textbf{Design} \\
          \midrule
          Energy threshold & 0.244 & 1.9 ms & Baseline \\
          MFCC + ZCR & 0.502 & 2.6 ms & Improved \\
          \rowcolor{lightgray}
          \textbf{Spectral (ours)} & \textbf{0.528} & \textbf{6.2 ms} & \textbf{Proposed} \\
          \bottomrule
        \end{tabular}
      \end{center}

      \vspace{0.3cm}
      \textbf{Avengers Trailer Results}:
      \begin{center}
        \small
        \begin{tabular}{@{}lccc@{}}
          \toprule
          Method & Segments & Speech (s) & Ratio \\
          \midrule
          Energy & 19 & 43.4 & 0.478 \\
          MFCC+ZCR & 20 & 70.6 & 0.650 \\
          \rowcolor{lightgray}
          \textbf{Spectral} & \textbf{13} & \textbf{78.5} & \textbf{0.700} \\
          \bottomrule
        \end{tabular}
      \end{center}
    \end{column}
    \begin{column}{0.42\textwidth}
      \begin{block}{Key Finding}
        Spectral VAD detects \textbf{2.17$\times$ more speech} than energy-only at only $3\times$ more compute (6.2 ms vs.\ 1.9 ms per utterance).
        \vspace{0.2cm}

        Critical for trailers where \textbf{40--70\% of audio is music-only}. Energy-based VAD fires during loud orchestral segments; spectral features discriminate harmonic music from voiced speech.
      \end{block}
      \begin{block}{Spectral VAD Discriminant}
        \vspace{-0.2cm}
        \[
          s = 0.45\,C_0 + 0.25\,\mathrm{cen} - 0.20\,\mathrm{flat} + 0.10\,(1-\mathrm{ZCR})
        \]
        \vspace{-0.3cm}
        \begin{itemize}
          \item $C_0$: MFCC energy
          \item cen: spectral centroid
          \item flat: spectral flatness (high $=$ noise-like)
          \item ZCR: zero-crossing rate
        \end{itemize}
      \end{block}
    \end{column}
  \end{columns}
\end{frame}

% ─────────────────────────────────────────────────────────────────────────────
\begin{frame}{Experiment 4: Enhancement Method Ablation}
  \begin{columns}[T]
    \begin{column}{0.55\textwidth}
      \textbf{Setup}: LibriSpeech + movie soundtrack noise at varying SNR.\\
      \textbf{Systems}: raw, Wiener, spectral subtraction, adaptive routing.

      \vspace{0.3cm}
      \begin{center}
        \begin{tabular}{@{}lccc@{}}
          \toprule
          \textbf{Method} & \textbf{WER} $\downarrow$ & \textbf{CER} $\downarrow$ & \textbf{RTF} \\
          \midrule
          \rowcolor{lightgray}
          \textbf{Raw (no enh.)} & \textbf{0.221} & \textbf{0.071} & 0.032 \\
          Spectral subtraction & 0.227 & 0.075 & \textbf{0.030} \\
          Adaptive routing & 0.367 & 0.168 & 0.041 \\
          Wiener filter & 0.499 & 0.270 & 0.050 \\
          \bottomrule
        \end{tabular}
      \end{center}
      \vspace{0.3cm}
      \begin{alertblock}{Surprising Result}
        Raw audio \textbf{outperforms all enhancement methods} on average. Wiener filter is 2$\times$ worse than raw (WER 0.499 vs.\ 0.221).
      \end{alertblock}
    \end{column}
    \begin{column}{0.42\textwidth}
      \begin{block}{Why Wiener Fails on Movies}
        Wiener filter over-smooths the spectrum, suppressing:
        \begin{itemize}
          \item High-frequency \textbf{fricatives} (s, f, sh) along with noise
          \item \textbf{Formant structure} — the ASR backbone
          \item Introduces \textbf{musical noise} artifacts at medium SNR
        \end{itemize}
        Modern int8-quantised Whisper is \textbf{intrinsically noise-robust}. Classical pre-processing confuses the decoder.
      \end{block}
      \begin{block}{When Enhancement Helps}
        Spectral subtraction only helps at \textbf{very low SNR} ($\leq 0$ dB):
        \begin{itemize}
          \item WER at $-5$ dB: raw 0.567, spectral sub \textbf{0.548}
          \item At $\geq 10$ dB: raw always wins
        \end{itemize}
        $\Rightarrow$ Adaptive routing is the correct design.
      \end{block}
    \end{column}
  \end{columns}
\end{frame}

% ─────────────────────────────────────────────────────────────────────────────
\begin{frame}{Experiment 2: Noise Robustness Benchmark — 960 Conditions}
  \begin{columns}[T]
    \begin{column}{0.52\textwidth}
      \textbf{Scale}: 10 speakers $\times$ 4 noise types $\times$ 6 SNR levels $\times$ 4 systems = \textbf{960 conditions}

      \vspace{0.2cm}
      \textbf{Noise types}: white Gaussian, pink (1/f), babble, movie soundtrack\\
      \textbf{SNR range}: $-5$ to $+20$ dB

      \vspace{0.2cm}
      \textbf{WER by SNR (all noise types, all systems):}
      \begin{center}
        \scriptsize
        \begin{tabular}{@{}lcccccc@{}}
          \toprule
          \textbf{System} & $\mathbf{-5}$ & \textbf{0} & \textbf{5} & \textbf{10} & \textbf{15} & \textbf{20} \\
          \midrule
          Raw & 0.567 & 0.378 & 0.242 & 0.199 & 0.194 & 0.208 \\
          Spectral sub & \textbf{0.548} & \textbf{0.378} & 0.263 & 0.212 & 0.207 & \textbf{0.200} \\
          Adaptive & 0.731 & 0.555 & 0.451 & 0.373 & 0.299 & 0.269 \\
          Wiener & 0.796 & 0.644 & 0.493 & 0.399 & 0.326 & 0.287 \\
          \bottomrule
        \end{tabular}
      \end{center}

      \vspace{0.2cm}
      \textbf{Mean across all conditions:}
      \begin{center}
        \small
        \begin{tabular}{@{}lccc@{}}
          \toprule
          System & WER & CER & RTF \\
          \midrule
          \rowcolor{lightgray}
          \textbf{Raw} & \textbf{0.298} & \textbf{0.131} & 0.043 \\
          Spectral sub & 0.301 & 0.129 & \textbf{0.035} \\
          Adaptive & 0.446 & 0.234 & 0.051 \\
          Wiener & 0.491 & 0.265 & 0.058 \\
          \bottomrule
        \end{tabular}
      \end{center}
    \end{column}
    \begin{column}{0.45\textwidth}
      \begin{block}{Key Findings}
        \begin{enumerate}
          \item \textbf{Raw beats all enhancement at SNR $\geq$ 10 dB} — modern Whisper handles typical movie SNR natively.
          \vspace{0.1cm}
          \item \textbf{Spectral sub only 0.9\% worse than raw} overall (0.301 vs.\ 0.298) and \textbf{wins at $-5$ dB}.
          \vspace{0.1cm}
          \item \textbf{Adaptive routing underperforms} because incorrect routing decisions (enhancing already-clean segments) dominate the error budget.
          \vspace{0.1cm}
          \item \textbf{Wiener is 65\% worse than raw} overall (0.491 vs.\ 0.298) — never recommended for movie audio.
        \end{enumerate}
      \end{block}
      \begin{alertblock}{Design Implication}
        Use raw audio as the default; apply spectral subtraction only when estimated SNR $< 0$ dB. Avoid Wiener filter entirely for cinematic content.
      \end{alertblock}
    \end{column}
  \end{columns}
\end{frame}

% ─────────────────────────────────────────────────────────────────────────────
\begin{frame}{Experiment 5: Model Size Comparison — tiny.en vs.\ medium.en}
  \begin{columns}[T]
    \begin{column}{0.52\textwidth}
      \textbf{Setup}: LibriSpeech test-clean + soundtrack noise across SNR levels.

      \vspace{0.3cm}
      \textbf{tiny.en vs.\ base.en (LibriSpeech):}
      \begin{center}
        \small
        \begin{tabular}{@{}lccccc@{}}
          \toprule
          Model & Params & WER & CER & Lat. & Load \\
          \midrule
          \textbf{base.en} & 74M & \textbf{0.172} & \textbf{0.040} & 367 ms & 17.6 s \\
          tiny.en & 39M & 0.186 & 0.050 & \textbf{262 ms} & \textbf{360 ms} \\
          \bottomrule
        \end{tabular}
      \end{center}

      \vspace{0.3cm}
      \textbf{tiny.en vs.\ medium.en (Trailer evaluation):}
      \begin{center}
        \small
        \begin{tabular}{@{}lcccc@{}}
          \toprule
          Model & WER & RTF & Load & Use case \\
          \midrule
          tiny.en & 0.875 & \textbf{0.052} & \textbf{360 ms} & Live captioning \\
          \rowcolor{lightgray}
          \textbf{medium.en} & \textbf{0.522} & 0.175 & 2.0 s & Batch offline \\
          \bottomrule
        \end{tabular}
      \end{center}

      \vspace{0.2cm}
      \begin{center}
        \textbf{medium.en reduces WER by 38\% relative} vs.\ tiny.en on trailers.
      \end{center}
    \end{column}
    \begin{column}{0.45\textwidth}
      \begin{block}{RTF — All Models Real-Time on CPU}
        \begin{center}
          \small
          \begin{tabular}{@{}lccc@{}}
            \toprule
            Model & Min & Mean & Max \\
            \midrule
            tiny.en & 0.013 & 0.043 & 0.548 \\
            base.en & 0.018 & 0.046 & 0.600 \\
            medium.en & 0.045 & 0.175 & 0.404 \\
            \bottomrule
          \end{tabular}
        \end{center}
        All $< 1.0$ = faster than real-time on Apple M1 CPU (int8 quantised).
      \end{block}
      \begin{block}{Recommendation}
        \begin{itemize}
          \item \textbf{Batch offline subtitling}: \texttt{medium.en}\\
            (best accuracy, still real-time)
          \item \textbf{Live on-the-fly captioning}: \texttt{tiny.en}\\
            (49$\times$ faster load, 1.3$\times$ faster inference, 360 ms cold start)
          \item \textbf{Balanced}: \texttt{base.en}\\
            (7.8\% WER gain over tiny at 1.4$\times$ latency)
        \end{itemize}
      \end{block}
    \end{column}
  \end{columns}
\end{frame}

% ─────────────────────────────────────────────────────────────────────────────
\begin{frame}{Experiment 6: Subtitle Quality \& Readability}
  \begin{columns}[T]
    \begin{column}{0.55\textwidth}
      \textbf{Metrics}: CPS violations, Timestamp MAE, WER under 3 noise scenarios.\\
      \textbf{Dataset}: LibriSpeech with synthetic noise.

      \vspace{0.3cm}
      \begin{center}
        \small
        \begin{tabular}{@{}lcccccc@{}}
          \toprule
          \textbf{Scenario} & \textbf{SNR} & \textbf{Cues} & \textbf{TS-MAE} & \textbf{WER} & \textbf{CPS Viol.} \\
          \midrule
          Clean (sndtrk) & 30 dB & 23 & 29.2 s & 0.966 & 7 \\
          Moderate (pink) & 10 dB & 21 & 24.2 s & 0.923 & 7 \\
          Heavy (sndtrk) & 0 dB & 16 & 22.3 s & 1.000 & 3 \\
          \bottomrule
        \end{tabular}
      \end{center}

      \vspace{0.2cm}
      \textbf{CPS Violations Analysis}:
      \begin{itemize}
        \item Violations occur when ASR emits long contiguous segments without word-boundary breaks
        \item \textbf{Not caused by transcription errors} — the text is correct but the cue is too long
        \item Max-words-per-cue rule ($\leq$12 words) reduces violations by $>$50\%
        \item Timestamp MAE improves at lower SNR because fewer cues are generated (shorter segments)
      \end{itemize}
    \end{column}
    \begin{column}{0.42\textwidth}
      \begin{block}{Netflix Compliance Targets}
        \begin{itemize}
          \item $\leq$ 20 chars/sec (CPS)
          \item $\leq$ 84 chars per cue (2 lines $\times$ 42 chars)
          \item $\geq$ 5 chars displayed for $\geq$ 1 second
          \item No subtitle overlap
        \end{itemize}
      \end{block}
      \begin{block}{Key Finding}
        CPS violations are a \textbf{post-processing concern}, not an ASR quality issue. The subtitle formatter (Stage 6) resolves them without any reprocessing by splitting long ASR outputs at word boundaries with duration proportioning.
      \end{block}
      \begin{alertblock}{TS-MAE Interpretation}
        High TS-MAE ($>$20 s) reflects the mismatch between ASR segment timestamps and the reference SRT timing, which was manually curated for cinematic effect. Not a fair metric for generated subtitles.
      \end{alertblock}
    \end{column}
  \end{columns}
\end{frame}

% ─────────────────────────────────────────────────────────────────────────────
\section{Trailer Evaluation}
\begin{frame}{Experiment 7: Mass Trailer Evaluation — 14 Trailers, 8 Genres}
  \textbf{Setup}: \texttt{medium.en} on 14 YouTube trailers with reference SRTs. Metrics: WER, CER, chrF, BLEU-1, RTF.

  \vspace{0.1cm}
  \begin{center}
    \scriptsize
    \begin{tabular}{@{}llccccc@{}}
      \toprule
      \textbf{Trailer} & \textbf{Genre} & \textbf{WER} $\downarrow$ & \textbf{CER} $\downarrow$ & \textbf{chrF} $\uparrow$ & \textbf{BLEU-1} $\uparrow$ & \textbf{RTF} \\
      \midrule
      Avengers Endgame \textit{(official SRT)} & action & \textbf{0.144} & 0.123 & \textbf{84.9} & \textbf{84.5} & 0.115 \\
      Spider-Man \textit{(official SRT)} & action & \textbf{0.183} & \textbf{0.090} & \textbf{86.3} & \textbf{87.4} & 0.190 \\
      Whiplash & drama & \textbf{0.405} & \textbf{0.266} & \textbf{68.9} & 67.0 & 0.201 \\
      Zodiac & thriller & \textbf{0.405} & \textbf{0.266} & \textbf{68.9} & 67.0 & 0.217 \\
      Bridesmaids & comedy & 0.449 & 0.325 & 61.3 & 64.3 & 0.254 \\
      Mad Max: Fury Road & action & 0.450 & 0.264 & 63.5 & 56.9 & 0.187 \\
      Dunkirk & action & 0.513 & 0.392 & 58.8 & 49.8 & \textbf{0.052} \\
      The Dark Knight & action & 0.525 & 0.355 & 61.2 & 54.4 & \textbf{0.045} \\
      Nomadland & drama & 0.533 & 0.393 & 63.7 & 62.7 & 0.404 \\
      Spider-Man: Spider-Verse & action & 0.570 & 0.514 & 48.8 & 36.4 & 0.081 \\
      The Revenant & drama & 0.554 & 0.566 & 47.4 & 32.6 & 0.287 \\
      The Social Network & drama & 0.634 & 0.488 & 57.1 & 47.0 & 0.204 \\
      Avengers: Age of Ultron \textit{(YT auto)} & action & 0.948 & 0.942 & 6.0 & 0.0 & 0.108 \\
      \midrule
      \textbf{Mean (all 14)} & — & 0.522 & 0.427 & 55.5 & 57.1 & 0.175 \\
      \bottomrule
    \end{tabular}
  \end{center}

  \begin{columns}[T]
    \begin{column}{0.5\textwidth}
      \textbf{Genre-level performance:}
      \begin{center}
        \scriptsize
        \begin{tabular}{@{}lccc@{}}
          \toprule
          Genre & WER & CER & chrF \\
          \midrule
          thriller & \textbf{0.405} & \textbf{0.266} & \textbf{68.9} \\
          comedy & 0.449 & 0.325 & 61.3 \\
          drama & 0.531 & 0.428 & 59.3 \\
          action & 0.593 & 0.513 & 46.2 \\
          \bottomrule
        \end{tabular}
      \end{center}
    \end{column}
    \begin{column}{0.48\textwidth}
      \small
      \textbf{Key}: Action trailers hardest (dense music, rapid cuts). Thriller/drama easier (more dialogue, cleaner audio). Official SRTs give dramatically better WER than YouTube auto-captions.
    \end{column}
  \end{columns}
\end{frame}

% ─────────────────────────────────────────────────────────────────────────────
\begin{frame}{The YouTube Ground Truth Fallacy}
  \begin{columns}[T]
    \begin{column}{0.52\textwidth}
      \textbf{Problem}: YouTube auto-generated SRTs are used as ``ground truth'' for trailer evaluation. But they are themselves produced by Google ASR.

      \vspace{0.3cm}
      \begin{center}
        \small
        \begin{tabular}{@{}lcc@{}}
          \toprule
          \textbf{Reference SRT Type} & \textbf{WER} & \textbf{chrF} \\
          \midrule
          Official human-authored & \textbf{0.144} & \textbf{84.9} \\
          Official human-authored & \textbf{0.183} & \textbf{86.3} \\
          YouTube auto-captions & 0.52–0.95 & 6–69 \\
          \bottomrule
        \end{tabular}
      \end{center}

      \vspace{0.3cm}
      \textbf{Avengers: Age of Ultron (YouTube auto-caption ref.)}
      \begin{itemize}
        \item WER = 0.948, chrF = 6.0
        \item The reference SRT contains \textbf{hallucinated lyrics} during orchestral sequences
        \item Our model correctly outputs silence; the reference outputs ``\textit{I don't want to talk about it...}''
        \item This inflates WER by penalising correct silence as substitution errors
      \end{itemize}
    \end{column}
    \begin{column}{0.45\textwidth}
      \begin{alertblock}{Key Insight}
        When evaluated against \textbf{official human-authored SRTs}:
        \begin{itemize}
          \item WER: \textbf{0.144} (Avengers Endgame)
          \item WER: \textbf{0.183} (Spider-Man)
          \item chrF: \textbf{84.9} and \textbf{86.3}
        \end{itemize}
        These match or exceed YouTube's own ASR on the same content.
      \end{alertblock}
      \begin{block}{Recommendation}
        \begin{enumerate}
          \item Always validate reference SRTs manually before computing WER
          \item Report chrF alongside WER (more robust to hallucination artifacts)
          \item Use Whisper confidence score to detect when reference may be wrong
        \end{enumerate}
      \end{block}
      \vspace{0.2cm}
      \textbf{Genre classifier accuracy}: \textbf{76.9\%} (leave-one-out k-NN on MFCC features) — useful for adaptive model selection per genre.
    \end{column}
  \end{columns}
\end{frame}

% ─────────────────────────────────────────────────────────────────────────────
\section{Novel Experiments}
\begin{frame}{Novel Experiment A: PESQ \& STOI Enhancement Quality}
  \begin{columns}[T]
    \begin{column}{0.55\textwidth}
      \textbf{Motivation}: WER is an end-to-end metric conflating ASR and pre-processing quality. PESQ and STOI measure purely the enhancement benefit, independent of the ASR backend.

      \vspace{0.3cm}
      \begin{block}{Metrics}
        \begin{itemize}
          \item \textbf{PESQ} (ITU-T P.862.2): Perceptual Evaluation of Speech Quality. Scale 0--4.5 MOS. Standard for telephony codec evaluation.
          \item \textbf{STOI} (Taal et al.\ 2011): Short-Time Objective Intelligibility. Scale 0--1. Correlates with human intelligibility scores.
        \end{itemize}
      \end{block}

      \vspace{0.2cm}
      \begin{center}
        \begin{tabular}{@{}lcc@{}}
          \toprule
          \textbf{Method} & \textbf{PESQ} $\uparrow$ & \textbf{STOI} $\uparrow$ \\
          \midrule
          Noisy (no enhancement) & 1.778 & 0.879 \\
          \rowcolor{lightgray}
          \textbf{Spectral subtraction} & \textbf{2.093} & \textbf{0.878} \\
          Wiener filter & 1.212 & 0.799 \\
          \bottomrule
        \end{tabular}
      \end{center}
    \end{column}
    \begin{column}{0.42\textwidth}
      \begin{block}{Key Findings}
        \begin{itemize}
          \item Spectral subtraction: \textbf{+17.8\% PESQ} over noisy baseline — perceptible quality improvement
          \item Wiener filter: \textbf{$-$31.8\% PESQ} — worse than no enhancement! Introduces musical noise artifacts
          \item STOI barely changes (0.879 $\to$ 0.878): intelligibility already high at typical movie SNR; enhancement is marginal
        \end{itemize}
      \end{block}
      \begin{alertblock}{Consistency with WER}
        PESQ and WER agree: spectral subtraction is best enhancement method; Wiener filter is harmful. This cross-validates our WER findings with an \textbf{independent perceptual quality metric}.
      \end{alertblock}
      \vspace{0.2cm}
      \small \textbf{Spectral subtraction parameters}: over-subtraction $\alpha=1.5$, noise estimated from leading frames, spectral floor = 5\% of clean magnitude.
    \end{column}
  \end{columns}
\end{frame}

% ─────────────────────────────────────────────────────────────────────────────
\begin{frame}{Novel Experiment B: Noise-Type Classifier}
  \begin{columns}[T]
    \begin{column}{0.52\textwidth}
      \textbf{Goal}: Automatically identify noise type in each audio segment to enable noise-specific enhancement routing.

      \vspace{0.2cm}
      \textbf{16-Dimensional Feature Vector}:
      \begin{itemize}
        \item MFCC $C_1$--$C_6$ (6): mean cepstral coefficients
        \item Log RMS (1): signal energy
        \item ZCR (1): zero-crossing rate
        \item Spectral centroid (1), flatness (1), rolloff (1)
        \item Sub-band energy ratios (4): $<$300 Hz, 300--1k, 1k--4k, $>$4k Hz
        \item Periodicity (1): normalised autocorrelation peak
      \end{itemize}

      \vspace{0.2cm}
      \textbf{Classes}: silence, white noise, pink noise, music/soundtrack, babble

      \begin{center}
        \begin{tabular}{@{}lc@{}}
          \toprule
          \textbf{Class} & \textbf{Accuracy} \\
          \midrule
          Silence & \textbf{1.00} \\
          White noise & 0.52 \\
          Music / soundtrack & 0.40 \\
          Pink noise & 0.00 \\
          Babble & 0.00 \\
          \midrule
          \textbf{Overall} & \textbf{0.316} \\
          \bottomrule
        \end{tabular}
      \end{center}
    \end{column}
    \begin{column}{0.45\textwidth}
      \begin{alertblock}{Failure Analysis}
        Pink noise and babble have similar MFCC profiles (both broad-spectrum, low periodicity). Rule-based features cannot discriminate them without temporal or higher-order statistics.
      \end{alertblock}
      \begin{block}{What the Classifier Gets Right}
        \begin{itemize}
          \item \textbf{Silence}: Perfect (100\%) — trivially separated by RMS
          \item \textbf{White noise}: Moderate (52\%) — unique flatness signature
          \item \textbf{Music}: Partial (40\%) — harmonic content partially captured by MFCC $C_1$--$C_6$
        \end{itemize}
      \end{block}
      \begin{block}{Next Step Motivated}
        A learned front-end classifier (MLP or SVM on mel-filterbank features) would separate pink/babble using temporal patterns. This is the clearest extension point for future work.
      \end{block}
    \end{column}
  \end{columns}
\end{frame}

% ─────────────────────────────────────────────────────────────────────────────
\begin{frame}{Novel Experiment C: Whisper Confidence Calibration}
  \begin{columns}[T]
    \begin{column}{0.52\textwidth}
      \textbf{Question}: Does Whisper's \texttt{avg\_logprob} reliably predict transcription quality?

      \vspace{0.2cm}
      \textbf{Method}: Correlate \texttt{avg\_logprob} with WER across \textbf{955 noise conditions} from Experiment 2. Compute calibration error and AUROC for binary error detection (WER $> 0.3$).

      \vspace{0.2cm}
      \begin{center}
        \begin{tabular}{@{}lc@{}}
          \toprule
          \textbf{Metric} & \textbf{Value} \\
          \midrule
          Pearson $r$ & $-0.783$ \\
          Spearman $\rho$ & $-0.859$ \\
          Expected Calibration Error (ECE) & 0.576 \\
          \rowcolor{lightgray}
          \textbf{AUROC (WER $>$ 0.3)} & \textbf{0.929} \\
          \bottomrule
        \end{tabular}
      \end{center}

      \vspace{0.2cm}
      \textbf{Interpretation}:
      \begin{itemize}
        \item Strong monotonic correlation (Spearman $-0.859$): as log-prob drops, WER rises
        \item Poor absolute calibration (ECE 0.576): log-prob values are not probabilistically meaningful
        \item But excellent \textbf{ranking/detection}: AUROC 0.929 means a threshold on \texttt{avg\_logprob} catches $\sim$93\% of bad segments
      \end{itemize}
    \end{column}
    \begin{column}{0.45\textwidth}
      \begin{block}{Practical Application}
        \textbf{Confidence-gated SRT filter}:
        \vspace{0.1cm}
        \begin{enumerate}
          \item Run ASR on each segment
          \item If \texttt{avg\_logprob} $< \tau$ (e.g.\ $-1.2$): flag as low-confidence
          \item Optionally: suppress the cue, mark it for human review, or trigger re-processing with a larger model
        \end{enumerate}
        \vspace{0.1cm}
        At AUROC = 0.929, a threshold set for 90\% recall of bad segments yields only 20\% false-alarm rate.
      \end{block}
      \begin{alertblock}{Already Implemented}
        Our pipeline uses \texttt{log\_prob\_threshold = $-$1.2} in faster-whisper to discard low-confidence segments before SRT output. This is the practical deployment of this finding.
      \end{alertblock}
    \end{column}
  \end{columns}
\end{frame}

% ─────────────────────────────────────────────────────────────────────────────
\begin{frame}{Novel Experiment D: Streaming Pipeline Latency}
  \begin{columns}[T]
    \begin{column}{0.52\textwidth}
      \textbf{Goal}: Simulate a \textbf{chunk-based live captioning} pipeline. Measure first-word latency and WER degradation as chunk size increases.

      \vspace{0.3cm}
      \begin{center}
        \begin{tabular}{@{}lccc@{}}
          \toprule
          \textbf{Mode} & \textbf{Latency} $\downarrow$ & \textbf{WER} $\downarrow$ & \textbf{RTF} \\
          \midrule
          Batch (full utt.) & 255 ms & \textbf{0.174} & \textbf{0.032} \\
          Stream 1 s chunks & \textbf{195 ms} & 0.607 & 0.260 \\
          Stream 2 s chunks & 186 ms & 0.324 & 0.232 \\
          \rowcolor{lightgray}
          \textbf{Stream 3 s chunks} & \textbf{203 ms} & \textbf{0.291} & 0.087 \\
          Stream 5 s chunks & 233 ms & 0.262 & 0.142 \\
          \bottomrule
        \end{tabular}
      \end{center}

      \vspace{0.2cm}
      \textbf{Why latency is non-monotonic}:
      \begin{itemize}
        \item Smaller chunks $\to$ less audio context $\to$ more VAD boundaries $\to$ processing overhead
        \item 3 s chunks hit a sweet spot between context and boundary overhead
      \end{itemize}
    \end{column}
    \begin{column}{0.45\textwidth}
      \begin{block}{Key Finding: 3-Second Chunks}
        \textbf{203 ms} first-word latency (well within 500 ms human perception threshold), \textbf{WER = 0.291} (only 67\% degradation vs.\ batch).

        This is the practical choice for \textbf{real-time live captioning} of movies and streaming content.
      \end{block}
      \begin{block}{The Batch vs.\ Stream Tradeoff}
        \begin{itemize}
          \item \textbf{Batch}: WER 0.174 but 255 ms latency — needs full utterance before transcription starts
          \item \textbf{3 s stream}: WER 0.291 but 203 ms — transcript appears while speaker is still talking
          \item 5 s chunks almost match batch WER (0.262) with 233 ms latency — viable for semi-live use
        \end{itemize}
      \end{block}
      \begin{alertblock}{1-second chunks fail}
        WER = 0.607: insufficient acoustic context for Whisper's attention mechanism to reliably decode words at sentence boundaries.
      \end{alertblock}
    \end{column}
  \end{columns}
\end{frame}

% ─────────────────────────────────────────────────────────────────────────────
\section{Scene Understanding}
\begin{frame}{Scene Understanding \& Mood Analysis}
  \begin{columns}[T]
    \begin{column}{0.52\textwidth}
      \textbf{Foote Novelty Score — Scene Boundary Detection}
      \begin{itemize}
        \item Extract per-frame features: MFCC$\times$13, chroma$\times$12, spectral contrast$\times$6, HPSS harmonic/percussive ratio, onset strength, centroid, rolloff, flatness, ZCR
        \item Compute \textbf{self-similarity matrix} of MFCC+chroma features
        \item Convolve with \textbf{Gaussian checkerboard kernel}
        \item Peaks = structural scene boundaries (dialogue$\to$action, music$\to$speech)
      \end{itemize}

      \vspace{0.3cm}
      \textbf{Mood Classification Labels}:
      \begin{center}
        \scriptsize
        \begin{tabular}{@{}lll@{}}
          \toprule
          \textbf{Mood} & \textbf{Acoustic Signature} & \textbf{Context} \\
          \midrule
          tense & High energy, low valence & Thriller/horror build-up \\
          epic & High energy, high harm. & Action/superhero \\
          sad & Low energy, low valence & Drama/tragedy \\
          happy & Med. energy, fast tempo & Comedy/uplifting \\
          calm & Low energy, high harm. & Quiet dialogue \\
          dialogue & Speech-like ZCR & Direct speech \\
          silence & RMS $< -40$ dB & Scene gap \\
          \bottomrule
        \end{tabular}
      \end{center}
    \end{column}
    \begin{column}{0.45\textwidth}
      \begin{block}{Multi-Modal Output}
        \textbf{Clean SRT}: Standard transcription for WER evaluation.

        \textbf{Annotated SRT}: In-line contextual tags:
        \begin{itemize}
          \item Moods: [TENSE]\,\textrm{😬}, [EPIC]\,\textrm{⚡}, [SAD]\,\textrm{😢}, [HAPPY]\,\textrm{😊}
          \item Events: [MUSIC]\,\textrm{🎵}, [IMPACT]\,\textrm{💥}, [SILENCE]
        \end{itemize}

        \textbf{Timeline JSON}: Frame-level acoustic properties (RMS, pitch, mood label) at 0.5 s resolution.

        \textbf{5-Panel Plot} per trailer:
        \begin{enumerate}
          \item Waveform + VAD segments
          \item Mel spectrogram
          \item Foote novelty score + boundaries
          \item Temporal mood timeline
          \item Mood distribution pie chart
        \end{enumerate}
      \end{block}
      \begin{alertblock}{Accessibility Value}
        Annotated SRT elevates subtitles from plain transcription to a \textbf{rich accessibility feature} for deaf/hard-of-hearing viewers, providing emotional context that hearing viewers get from the soundtrack.
      \end{alertblock}
    \end{column}
  \end{columns}
\end{frame}

% ─────────────────────────────────────────────────────────────────────────────
\section{Conclusions}
\begin{frame}{Key Findings Summary}
  \begin{columns}[T]
    \begin{column}{0.5\textwidth}
      \textbf{Core findings:}
      \begin{enumerate}
        \item \textbf{Raw ASR beats enhancement at SNR $\geq$ 10 dB.}\\
          Modern int8 Whisper is intrinsically noise-robust at typical movie SNRs.

        \item \textbf{Adaptive routing is the right design.}\\
          Withhold enhancement for clean segments; apply spectral subtraction only below 0 dB.

        \item \textbf{Spectral subtraction $>$ Wiener} for movie audio.\\
          PESQ +17.8\%, WER $-38\%$ relative vs.\ Wiener under soundtrack noise.

        \item \textbf{Spectral VAD detects 2.17$\times$ more speech} than energy-only at only 3$\times$ extra compute (6.2 ms vs.\ 1.9 ms).

        \item \textbf{medium.en reduces WER by 38\%} vs.\ tiny.en on trailers (0.522 vs.\ 0.875), still real-time (RTF 0.175).
      \end{enumerate}
    \end{column}
    \begin{column}{0.5\textwidth}
      \textbf{Novel findings:}
      \begin{enumerate}
        \setcounter{enumi}{5}
        \item \textbf{Whisper avg\_logprob is a strong error predictor.}\\
          Spearman $\rho = -0.859$, AUROC = 0.929 for WER $>$ 0.3 detection.

        \item \textbf{3-second chunks} give best streaming tradeoff.\\
          203 ms latency, WER = 0.291 (vs.\ 255 ms / 0.174 for batch).

        \item \textbf{Rule-based noise classifier fails} on overlapping classes (31.6\% accuracy). Motivates a learned front-end.

        \item \textbf{YouTube auto-captions are flawed ground truth.}\\
          Official SRT WER = 0.144 (Avengers) vs.\ 0.948 with YouTube auto-captions.

        \item \textbf{Foote novelty score} successfully detects structural scene transitions for annotated subtitle generation.
      \end{enumerate}
    \end{column}
  \end{columns}
\end{frame}

% ─────────────────────────────────────────────────────────────────────────────
\begin{frame}{Future Work \& Conclusions}
  \begin{columns}[T]
    \begin{column}{0.52\textwidth}
      \textbf{Immediate extensions:}
      \begin{itemize}
        \item \textbf{Learned noise classifier}: Replace rule-based 16-dim classifier with MLP/SVM on mel-filterbank features to separate pink noise from babble (current accuracy: 31.6\%)
        \item \textbf{LLM post-processing}: Use GPT-4o or Claude to correct ASR errors using dialogue context and movie script priors
        \item \textbf{GPU deployment}: On NVIDIA A10, medium.en RTF would drop from 0.175 to $\sim$0.02 (10$\times$ speedup)
        \item \textbf{Speaker diarisation}: Attribute subtitle cues to specific characters using pyannote.audio
        \item \textbf{Confidence-gated re-run}: Automatically re-run low-confidence segments with medium.en when tiny.en is used in live mode
      \end{itemize}
    \end{column}
    \begin{column}{0.45\textwidth}
      \begin{block}{What We Built}
        A \textbf{research-grade auto-subtitling system} combining:
        \begin{itemize}
          \item Classical DSP (VAD, spectral subtraction, Foote novelty)
          \item Neural ASR (faster-whisper, int8, Silero VAD)
          \item Multi-dimensional evaluation (WER, CER, PESQ, STOI, chrF, CPS, AUROC)
          \item 11 experiments across 960+ conditions
          \item 14 trailers $\times$ 8 genres
          \item 25+ figures, 20+ result tables
        \end{itemize}
      \end{block}
      \begin{alertblock}{The Main Takeaway}
        \textbf{Do not blindly apply classical enhancement to neural ASR.} Whisper's internal noise robustness makes spectral pre-processing harmful for SNR $\geq$ 10 dB. Adaptive, confidence-aware routing is the correct engineering approach for production movie subtitling.
      \end{alertblock}
    \end{column}
  \end{columns}
  \vspace{0.4cm}
  \begin{center}
    \large\textbf{Thank you!} \quad Code: \texttt{/project/*.py} \quad Demo: \texttt{python run\_enhanced.py}
  \end{center}
\end{frame}

% ─────────────────────────────────────────────────────────────────────────────
\appendix
\begin{frame}{Appendix: All Experiments at a Glance}
  \begin{center}
    \small
    \begin{tabular}{@{}clll@{}}
      \toprule
      \textbf{\#} & \textbf{Name} & \textbf{Dataset} & \textbf{Key Result} \\
      \midrule
      1 & Avengers Trailer End-to-End & Avengers trailer (125 s) & RTF = 0.048, speech 70\% \\
      2 & Noise Robustness Benchmark & LibriSpeech × 960 conditions & Raw WER = 0.298 (best) \\
      3 & VAD Ablation & LibriSpeech + trailer & Spectral: 2.17× more speech \\
      4 & Enhancement Ablation & LibriSpeech + soundtrack & Wiener 2× worse than raw \\
      5 & Model Size Comparison & LibriSpeech & base.en WER 0.172, load 17.6 s \\
      6 & Subtitle Quality & LibriSpeech + 3 scenarios & CPS fix: $\leq$12 words/cue \\
      7 & Mass Trailer Eval & 14 trailers, 8 genres & Mean WER 0.522 (medium.en) \\
      8 (A) & PESQ \& STOI & Synthetic noise & Spectral sub: PESQ +17.8\% \\
      9 (B) & Noise Classifier & Synthetic noise, 5 classes & Overall accuracy: 31.6\% \\
      10 (C) & Confidence Calibration & 955 conditions (Exp 2) & AUROC = 0.929 \\
      11 (D) & Streaming Latency & LibriSpeech, 5 chunk sizes & 3 s: 203 ms latency \\
      \bottomrule
    \end{tabular}
  \end{center}
\end{frame}

\begin{frame}{Appendix: Experiment 1 — Avengers Trailer Deep Dive}
  \begin{columns}[T]
    \begin{column}{0.52\textwidth}
      \textbf{Setup}: Full pipeline on 125-second Avengers Endgame trailer.\\
      \textbf{Script}: \texttt{run\_full\_experiments.py} $\to$ \texttt{exp1\_trailer()}

      \vspace{0.3cm}
      \begin{center}
        \begin{tabular}{@{}lccccc@{}}
          \toprule
          \textbf{VAD} & \textbf{Segs} & \textbf{Speech (s)} & \textbf{Ratio} & \textbf{ms/utt} \\
          \midrule
          Energy & 19 & 43.4 & 0.478 & 1.9 \\
          MFCC+ZCR & 20 & 70.6 & 0.650 & 2.6 \\
          \textbf{Spectral} & \textbf{13} & \textbf{78.5} & \textbf{0.700} & \textbf{6.2} \\
          \bottomrule
        \end{tabular}
      \end{center}

      \vspace{0.3cm}
      \textbf{ASR performance}:
      \begin{itemize}
        \item Overall RTF: \textbf{0.048} (20$\times$ faster than real-time)
        \item Subtitle cues generated: \textbf{9}
        \item CPS violations: \textbf{2}
        \item Enhancement triggered: 4 of 13 segments (SNR $< 18$ dB)
      \end{itemize}
    \end{column}
    \begin{column}{0.45\textwidth}
      \begin{block}{Generated Outputs}
        \begin{itemize}
          \item \texttt{predicted.srt}: Clean SRT (9 cues)
          \item \texttt{predicted\_annotated.srt}: With [EPIC], [TENSE], [MUSIC] tags
          \item \texttt{scene\_analysis.png}: 5-panel plot
          \item \texttt{scene\_analysis.json}: Foote boundaries + mood timeline
          \item \texttt{metrics.json}: WER, CER, RTF, CPS
        \end{itemize}
      \end{block}
      \begin{block}{Dominant Mood Distribution}
        \begin{itemize}
          \item Epic: 68\% of audio duration
          \item Tense: 18\%
          \item Dialogue: 10\%
          \item Silence: 4\%
        \end{itemize}
        $\Rightarrow$ Consistent with action blockbuster genre.
      \end{block}
    \end{column}
  \end{columns}
\end{frame}

\end{document}
